\begin{lemma}[Zero-energy BNT vectors lie in the parent-Hamiltonian kernel]
    \label{lem:bnt_span_le_parent_hamiltonian_kernel}
    \lean{MPSTensor.bntMPSVectorSpan_le_ker_parentHamiltonian_of_forall_annihilates}
    \lean{MPSTensor.bntMPSVectorSpan_le_ker_parentHamiltonian_of_forall_frustrationFree}
    \lean{MPSTensor.bntMPSVectorSpan_le_ker_parentHamiltonian_toTensorFromBlocks}
    \leanok
    \uses{def:bnt_span, def:is_frustration_free, def:parent_hamiltonian,
        thm:chain_ground_space_implies_frustration_free}
    Let $B$ be a canonical-form tensor with BNT components
    $A_1,\ldots,A_g$. If every component vector has zero energy for the
    parent Hamiltonian, $H_L^{(N)}(B)V^{(N)}(A_j)=0$ for $1\le j\le g$, then
    \begin{align}
        \spn\{V^{(N)}(A_j):j=1,\ldots,g\}
        &\subseteq
        \ker H_L^{(N)}(B). \notag
    \end{align}
    It is enough to assume the local frustration-free equations
    $h_iV^{(N)}(A_j)=0$ for every $i$ and $j$.
    For $N>0$ and $L\le N$, in the block-diagonal case
    $B=\bigoplus_j\mu_jA_j$, with every $\mu_j\ne0$, the same inclusion
    follows because each local space $G_L(A_j)$ is contained in $G_L(B)$.
    This is the easy inclusion in the source ground-space spanning equation of
    \cite[Definition~3.9]{Cirac2016MPDO_arXiv}; the reverse inclusion is the
    remaining spanning direction.
\end{lemma}

\begin{proof}
    \leanok
    The kernel is a linear subspace. Hence it is enough to check the generators
    $V^{(N)}(A_j)$. The first hypothesis gives this directly. Under the
    local frustration-free equations,
    \begin{align}
        H_L^{(N)}(B)V^{(N)}(A_j)
        &=\sum_i h_iV^{(N)}(A_j) =0. \notag
    \end{align}
    For a block-diagonal tensor, write $R_{i,\tau}$ for the cyclic
    length-$L$ restriction at site $i$, with the outside configuration
    fixed by $\tau$. Then
    $R_{i,\tau}V^{(N)}(A_j) \in G_L(A_j)\subseteq G_L(B)$.
    The preceding theorem therefore gives
    $h_i(B,L)V^{(N)}(A_j)=0$ for every $i$.
\end{proof}

\begin{lemma}[Kernel containment from chain-space splitting]
    \label{lem:block_diagonal_kernel_to_bnt_span_from_chain_splitting}
    \lean{MPSTensor.ker_parentHamiltonian_toTensorFromBlocks_le_bntMPSVectorSpan}
    \leanok
    \uses{thm:parent_hamiltonian_kernel_le_chain_ground_space,
        def:parent_hamiltonian, def:chain_ground_space, def:bnt_span,
        def:mpv_submodule, lem:bnt_span_eq_sum_mpv_submodules}
    Let $B=\bigoplus_{j=0}^{r-1}\mu_jA_j$. Suppose $N>0$, $L\le N$, and
    the periodic chain constraints for $B$ split into the block constraints:
    \begin{align}
        \mathcal G_{N,L}(B)
        &\subseteq
        \bigvee_{j=0}^{r-1}\mathcal G_{N,L}(A_j). \notag
    \end{align}
    If every block chain space is contained in its periodic MPS line,
    \begin{align}
        \mathcal G_{N,L}(A_j)
        &\subseteq \C V^{(N)}(A_j)
        \qquad (0\le j<r), \notag
    \end{align}
    then the zero-energy space of the parent Hamiltonian is contained in the
    BNT vector span:
    \begin{align}
        \ker H_L^{(N)}(B)
        &\subseteq
        \spn\{V^{(N)}(A_j):j=0,\ldots,r-1\}. \notag
    \end{align}
\end{lemma}

\begin{proof}
    \leanok
    Theorem~\ref{thm:parent_hamiltonian_kernel_le_chain_ground_space} gives
    $\ker H_L^{(N)}(B)\subseteq\mathcal G_{N,L}(B)$. Together with the
    splitting hypothesis,
    \begin{align}
        \ker H_L^{(N)}(B)
        &\subseteq \mathcal G_{N,L}(B) \subseteq \bigvee_j\mathcal G_{N,L}(A_j). \notag
    \end{align}
    The blockwise containments give
    \begin{align}
        \bigvee_j\mathcal G_{N,L}(A_j)
        &\subseteq \bigvee_j\C V^{(N)}(A_j), \notag
    \end{align}
    and Lemma~\ref{lem:bnt_span_eq_sum_mpv_submodules} gives
    \begin{align}
        \bigvee_j\C V^{(N)}(A_j)
        &=\spn\{V^{(N)}(A_j):j=0,\ldots,r-1\}. \notag
    \end{align}
\end{proof}

\begin{lemma}[Kernel containment from the block-diagonal boundary comparison]
    \label{lem:block_diagonal_kernel_to_bnt_span_from_pgvwc_comparison}
    \lean{MPSTensor.ker_parentHamiltonian_toTensorFromBlocks_le_bntMPSVectorSpan_of_pgvwc_comparison}
    \leanok
    \uses{lem:block_diagonal_kernel_to_bnt_span_from_chain_splitting,
        lem:bnt_block_diagonal_pgvwc_comparison_ground_space}
    Let $B=\bigoplus_{j=0}^{r-1}\mu_jA_j$, with every $\mu_j\ne0$.
    Suppose the physical alphabet and all block dimensions are nonempty. Assume
    that the block family is irreducible, left-canonical, has normalized
    self-overlap, and has no two blocks gauge-phase equivalent. Suppose also
    that each $A_j$ is $L_0$-block-injective, that $L_0>0$, and that
    $\sum_a A^j_a(A^j_a)^\dagger =I \qquad (0\le j<r)$.
    Let $N>0$, $L>0$, $L\le N$, $L+L_0\le N$, and
    $(L_0+1) +(r-1)((L_0+1)+((L_0+1)+(L_0+1)))+1 \le L$.
    Assume the opened-boundary $C^j,D^j$ comparison: whenever
    $\psi\in\mathcal G_{N,L}(B)$ and
    $\psi=\Gamma_N^B(\bigoplus_j X_j)$, there are matrices
    $C^j_{i,\rho}$ such that, for every boundary-crossing interval $i$,
    every $\rho$, and every $\beta$,
    $A^j_\beta C^j_{i,\rho} =((\mu_j^N X_j)A^j_\beta)A^j_\rho$.
    If every block chain space is contained in its periodic MPS line,
    \begin{align}
        \mathcal G_{N,L}(A_j)
        &\subseteq \C V^{(N)}(A_j)
        \qquad (0\le j<r), \notag
    \end{align}
    then
    \begin{align}
        \ker H_L^{(N)}(B)
        &\subseteq
        \spn\{V^{(N)}(A_j):j=0,\ldots,r-1\}. \notag
    \end{align}
\end{lemma}

\begin{proof}
    \leanok
    Lemma~\ref{lem:bnt_block_diagonal_pgvwc_comparison_ground_space} gives
    $\mathcal G_{N,L}(B) =\bigvee_{j=0}^{r-1}\mathcal G_{N,L}(A_j)$
    from the stated block hypotheses and the opened-boundary comparison.
    Lemma~\ref{lem:block_diagonal_kernel_to_bnt_span_from_chain_splitting}
    then applies to this equality and to the blockwise containments.
\end{proof}

\begin{lemma}[Kernel containment from boundary-crossing comparisons]
    \label{lem:block_diagonal_kernel_to_bnt_span_from_crossing_pgvwc_comparison}
    \lean{MPSTensor.ker_parentHamiltonian_toTensorFromBlocks_le_bntMPSVectorSpan_of_crossing_pgvwc_comparison}
    \leanok
    \uses{lem:block_diagonal_kernel_to_bnt_span_from_chain_splitting,
        lem:bnt_block_diagonal_crossing_pgvwc_comparison_ground_space}
    Let $B=\bigoplus_{j=0}^{r-1}\mu_jA_j$, with every $\mu_j\ne0$.
    Suppose the physical alphabet and all block dimensions are nonempty. Assume
    that the block family is irreducible, left-canonical, has normalized
    self-overlap, and has no two blocks gauge-phase equivalent. Suppose also
    that each $A_j$ is $L_0$-block-injective, that $L_0>0$, and that
    $\sum_a A^j_a(A^j_a)^\dagger =I \qquad (0\le j<r)$.
    Let $N>0$, $L>0$, $L\le N$, $L+L_0\le N$, and
    $(L_0+1) +(r-1)((L_0+1)+((L_0+1)+(L_0+1)))+1 \le L$.
    Assume that, for every boundary-crossing interval beginning at $i$, the
    simultaneous products
    $w\longmapsto(A^1_w,\ldots,A^r_w), \qquad |w|=N-i$,
    span the product algebra $\prod_j\MN{D_j}$. If every block chain space is
    contained in its periodic MPS line,
    \begin{align}
        \mathcal G_{N,L}(A_j)
        &\subseteq \C V^{(N)}(A_j)
        \qquad (0\le j<r), \notag
    \end{align}
    then
    \begin{align}
        \ker H_L^{(N)}(B)
        &\subseteq
        \spn\{V^{(N)}(A_j):j=0,\ldots,r-1\}. \notag
    \end{align}
\end{lemma}

\begin{proof}
    \leanok
    Lemma~\ref{lem:bnt_block_diagonal_crossing_pgvwc_comparison_ground_space}
    gives
    $\mathcal G_{N,L}(B) =\bigvee_{j=0}^{r-1}\mathcal G_{N,L}(A_j)$
    from the block hypotheses, the local boundary-crossing constraints, and the
    stated tail-word span.
    Lemma~\ref{lem:block_diagonal_kernel_to_bnt_span_from_chain_splitting}
    then applies to this equality and to the blockwise containments.
\end{proof}

\begin{lemma}[Block-diagonal spanning is equivalent to the reverse inclusion]
    \label{lem:block_diagonal_parent_spanning_reverse_inclusion}
    \lean{MPSTensor.hasParentHamiltonianGroundSpaceSpanning_toTensorFromBlocks_iff_ker_le_bntMPSVectorSpan}
    \lean{MPSTensor.hasNNCPHGroundSpaces_toTensorFromBlocks_iff_forall_isNNCPH_and_ker_le_bntMPSVectorSpan}
    \leanok
    \uses{def:parent_hamiltonian_ground_space_spanning,
        lem:bnt_span_le_parent_hamiltonian_kernel, def:has_nncph_ground_spaces}
    Let $B=\bigoplus_{j=0}^{r-1}\mu_jA_j$, with every $\mu_j\ne0$. Then, for
    every $N>L$, the source parent-Hamiltonian spanning condition
    $\ker H_L^{(N)}(B) =\spn\{V^{(N)}(A_j):j=0,\ldots,r-1\}$
    is equivalent to the reverse containment, for the same $N$:
    \begin{align}
        \ker H_L^{(N)}(B)
        &\subseteq
        \spn\{V^{(N)}(A_j):j=0,\ldots,r-1\}. \notag
    \end{align}
    In the nearest-neighbor case $L=2$, the all-chain source condition of
    Definition~\ref{def:has_nncph_ground_spaces} is equivalent to the
    conjunction
    \begin{align}
        &\left(\forall N>2,\ \forall i,j\in\{0,\ldots,N-1\},
        h_i(B,2)h_j(B,2)=h_j(B,2)h_i(B,2)\right) \notag\\
        &{}\land
        \left(\forall N>2,
        \ker H_2^{(N)}(B)
        \subseteq
        \spn\{V^{(N)}(A_j):j=0,\ldots,r-1\}\right). \notag
    \end{align}
\end{lemma}

\begin{proof}
    \leanok
    The forward direction of the displayed equality gives the containment by
    rewriting the kernel. Conversely, for $N>L$ the easy inclusion above gives
    \begin{align}
        \spn\{V^{(N)}(A_j):j=0,\ldots,r-1\}
        &\subseteq
        \ker H_L^{(N)}(B), \notag
    \end{align}
    and the assumed reverse containment gives equality. For $L=2$, the
    all-chain nearest-neighbor condition of
    Definition~\ref{def:has_nncph_ground_spaces} decomposes as the conjunction
    \begin{align}
        &\left(\forall N>2,\ \forall i,j\in\{0,\ldots,N-1\},
        h_i(B,2)h_j(B,2)=h_j(B,2)h_i(B,2)\right) \notag\\
        &{}\land
        \left(\forall N>2,
        \ker H_2^{(N)}(B)
        \subseteq
        \spn\{V^{(N)}(A_j):j=0,\ldots,r-1\}\right). \notag
    \end{align}
    The second conjunct is the reverse containment just identified with the
    source spanning clause.
\end{proof}

\begin{lemma}[Chain-space comparison gives the block-diagonal spanning equation]
    \label{lem:block_diagonal_chain_space_comparison_parent_spanning}
    \lean{MPSTensor.hasParentHamiltonianGroundSpaceSpanning_toTensorFromBlocks_of_chain_eq_iSup_mpv}
    \lean{MPSTensor.hasParentHamiltonianGroundSpaceSpanning_toTensorFromBlocks_of_chain_eq_iSup_chain}
    \leanok
    \uses{def:parent_hamiltonian_ground_space_spanning,
        thm:parent_hamiltonian_kernel_eq_chain_ground_space,
        lem:block_diagonal_parent_spanning_reverse_inclusion,
        lem:bnt_span_eq_sum_mpv_submodules}
    Let $B=\bigoplus_{j=0}^{r-1}\mu_jA_j$, where $\mu_j\ne0$. Suppose that,
    for every $N>L$,
    $G_L^{(N)}(B) =\bigvee_{j=0}^{r-1}\spn\{V^{(N)}(A_j)\}$.
    Then the source parent-Hamiltonian spanning equation holds:
    $\ker H_N(B,L) =\spn\{V^{(N)}(A_j):j=0,\ldots,r-1\}$.
    It is enough to prove the two equations
    \begin{align}
        G_L^{(N)}(B)
        &=\bigvee_{j=0}^{r-1}G_L^{(N)}(A_j), \notag\\
        G_L^{(N)}(A_j)
        &=\spn\{V^{(N)}(A_j)\}. \notag
    \end{align}
\end{lemma}

\begin{proof}
    \leanok
    The finite-volume identity gives
    $\ker H_N(B,L)=G_L^{(N)}(B)$. The assumed chain-space comparison therefore
    rewrites the kernel as
    $\bigvee_{j=0}^{r-1}\spn\{V^{(N)}(A_j)\}$. The latter supremum is the BNT
    vector span. Conversely, the easy inclusion of
    Lemma~\ref{lem:bnt_span_le_parent_hamiltonian_kernel} gives
    $\spn\{V^{(N)}(A_j):j=0,\ldots,r-1\} \subseteq\ker H_N(B,L)$,
    so the displayed equation is the source spanning clause. If the comparison
    is first stated as $G_L^{(N)}(B)=\bigvee_jG_L^{(N)}(A_j)$, the single-block
    equations $G_L^{(N)}(A_j)=\spn\{V^{(N)}(A_j)\}$ give the previous
    hypothesis.
\end{proof}

\begin{lemma}[Gauge invariance of fixed points and parent interactions]
    \label{lem:rfp_parent_gauge_invariance}
    \lean{MPSTensor.GaugeEquiv.hasPhysicalBlockingIsometry}
    \lean{MPSTensor.GaugeEquiv.hasPhysicalBlockingIsometry_iff}
    \lean{MPSTensor.GaugeEquiv.isTransferIdempotent_iff}
    \lean{MPSTensor.GaugeEquiv.parentInteraction_eq}
    \lean{MPSTensor.GaugeEquiv.localTerm_eq}
    \lean{MPSTensor.GaugeEquiv.isNNCPH}
    \leanok
    \uses{def:gauge_equiv, def:mps_transfer_idempotence,
        def:parent_interaction}
    A virtual similarity preserves the physical blocking relation and hence
    transfer idempotence. It also preserves every local MPS space, so the
    canonical parent interaction and all its translated local terms are
    unchanged.
\end{lemma}

\begin{proof}\leanok
    \uses{thm:physical_blocking_isometry_iff_transfer_idempotence}
    Suppose $B^i=XA^iX^{-1}$ for every physical letter. If the blocking
    relation for $A$ is
    $A^{i_1}A^{i_2} =\sum_j U^{j}_{i_1,i_2}A^j$,
    then conjugation by $X$ gives the same relation for $B$. Hence transfer
    idempotence is invariant under the similarity. At every length $L$, the
    two tensors have the same local MPS space, so $G_L(A)=G_L(B)$ and
    $h_i^{(N)}(A)=h_i^{(N)}(B)$. The parent-interaction identities and the
    commutation equations follow.
\end{proof}

\begin{theorem}[Renormalization fixed points have commuting nearest-neighbor local terms]\label{thm:rfp_implies_nncph}
    \lean{MPSTensor.rfp_implies_nncph}
    \leanok
    \uses{def:mps_transfer_idempotence, def:is_nncph, def:nt_cpsv}
    If $A$ has nonzero virtual dimension, is a normal tensor in the sense of
    \cite{Cirac2016MPDO_arXiv}, and is a renormalization fixed point, then for
    every chain length $N>2$ the two-site local terms commute. This theorem
    records the commutation part of the nearest-neighbor source condition, not
    the ground-space spanning part.
\end{theorem}

\begin{proof}\leanok
    \uses{thm:cpsv_normal_tp_gauge, lem:rfp_parent_gauge_invariance,
        thm:left_canonical_rfp_chain_commutation, thm:gauge_invariance_normal}
    Put $A$ into its trace-preserving gauge. This is a pure similarity, since
    the normal-tensor normalization fixes the transfer spectral radius at one.
    Similarity preserves the physical blocking relation and the local MPS
    spaces, hence also the canonical parent interactions. In the
    trace-preserving gauge, the Appendix~B virtual-bond projectors give
    $h_0^{(3)}h_1^{(3)} =h_1^{(3)}h_0^{(3)}$.
    Cyclic transport gives, for every $N>2$ and every site $i$,
    $h_i^{(N)}h_{i+1}^{(N)} =h_{i+1}^{(N)}h_i^{(N)}$.
    For every pair $i,j$ whose two-site supports are disjoint, locality gives
    $h_i^{(N)}h_j^{(N)} =h_j^{(N)}h_i^{(N)}$.
    These cases exhaust all pairs of two-site terms on the periodic chain.
    Transporting the parent interactions back proves the claim.
\end{proof}

\begin{theorem}[Renormalization fixed points satisfy the NNCPH commutation and zero-energy equations]%
    \label{thm:rfp_implies_nncph_ground_state}
    \lean{MPSTensor.rfp_implies_nncph_ground_state}
    \leanok
    \uses{thm:rfp_implies_nncph, def:is_nncph_ground_state}
    If $A$ has nonzero virtual dimension, is a normal tensor in the sense of
    \cite{Cirac2016MPDO_arXiv}, and is a renormalization fixed point, then for
    every chain length $N>2$ the MPS vector $V^{(N)}(A)$ satisfies the
    commutation and zero-energy equations for nearest-neighbor local terms:
    \begin{align}
        h_i h_j&=h_j h_i, \notag\\
        h_i V^{(N)}(A)&=0. \notag
    \end{align}
    The zero-energy equation is the usual parent-Hamiltonian
    frustration-freeness statement. The separate ground-space spanning clause
    of the source definition is not asserted here.
\end{theorem}

\begin{proof}
    \leanok
    \uses{thm:rfp_implies_nncph, lem:parent_hamiltonian_ff}
    The first equation is Theorem~\ref{thm:rfp_implies_nncph}. The second is
    Lemma~\ref{lem:parent_hamiltonian_ff} with $L=2$.
\end{proof}

\begin{theorem}[Nearest-neighbor parent ground spaces for a single normal RFP tensor]%
    \label{thm:rfp_single_normal_nncph_ground_spaces}
    \lean{MPSTensor.rfp_hasParentHamiltonianGroundSpaceSpanning_single}
    \lean{MPSTensor.rfp_implies_hasNNCPHGroundSpaces_single}
    \leanok
    \uses{def:mps_transfer_idempotence, def:nt_cpsv,
        def:parent_hamiltonian_ground_space_spanning,
        def:has_nncph_ground_spaces}
    Let $A$ have positive bond dimension and be a normal renormalization
    fixed-point tensor in the sense of \cite{Cirac2016MPDO_arXiv}. If the
    canonical-form tensor consists of the single normal sector $A$, then for
    every $N>2$,
    $\ker H_2^{(N)}(A) =\spn\bigl\{V^{(N)}(A)\bigr\}$.
    Its translated two-site interactions also commute and annihilate
    $V^{(N)}(A)$. Hence the singleton family satisfies the all-chain
    nearest-neighbor commuting parent-Hamiltonian ground-space condition.

    This statement is restricted to one normal sector. It does not assert the
    corresponding result for canonical-form tensors with several sectors or
    repeated copies.
\end{theorem}

\begin{proof}\leanok
    \uses{thm:cpsv_normal_eventually_injective, thm:rfp_nt_injective,
        thm:parent_hamiltonian_kernel_eq_chain_ground_space,
        thm:chain_ground_space_eq_mpv_blk,
        lem:bnt_span_eq_sum_mpv_submodules,
        thm:rfp_implies_nncph_ground_state}
    Spectral normality gives eventual block injectivity. The fixed-point
    equation then implies that the one-site matrices already span the full
    matrix algebra. Thus $A$ is injective. For $N>2$, the finite-volume kernel
    identity and injective periodic-chain theorem give
    \begin{align}
        \ker H_2^{(N)}(A)
        &=\mathcal G_{N,2}(A) =\spn\bigl\{V^{(N)}(A)\bigr\}. \notag
    \end{align}
    The span of the matrix product vectors of the singleton family is the last
    space in this equality. Theorem~\ref{thm:rfp_implies_nncph_ground_state}
    supplies commutation and zero energy on every chain of length greater than
    two.
\end{proof}

\begin{theorem}[Commuting nearest-neighbor terms for distinct multiplicity-one RFP sectors]%
    \label{thm:rfp_multiplicity_one_bnt_nncph}
    \lean{MPSTensor.IsBNTCanonicalForm.rfp_implies_nncph_basisDirectSum}
    \leanok
    \uses{def:sector_bnt_cf, def:mps_transfer_idempotence, def:is_nncph,
        lem:nncph_from_adjacent_two_site_commutation}
    Let $A_1,\ldots,A_g$ be the distinct basis tensors of a BNT canonical
    form, and let $B=\bigoplus_{j=1}^g A_j$ contain one copy of every basis
    tensor. If $B$ is a renormalization fixed point, then for every $N>2$
    its translated two-site parent interactions commute pairwise. Thus $B$
    satisfies the nearest-neighbor commuting parent-Hamiltonian condition.

    This statement treats one copy of every distinct sector. It does not
    include repeated copies or arbitrary raw sector weights.
\end{theorem}

\begin{proof}\leanok
    The fixed-point equation gives residual tensors $U_j$ whose matrix-entry
    columns form one orthonormal family on the disjoint union of the sector
    matrix spaces. On the tensor square of this disjoint union, let $Q$ be the
    sum of the sector bond projectors, extended by zero between unequal
    sectors. Write $L_j$ and $R_j$ for the left and right overlapping lifts
    of the sector-$j$ bond projector. The single-sector relation and the
    mixed-sector orthogonality give
    \begin{align}
        L_jR_j&=R_jL_j, \notag\\
        L_jR_k&=0=R_jL_k
        \qquad (j\ne k), \notag
    \end{align}
    where the latter identities follow from the shared one-site factor
    $E_j^*E_k=0$ for the canonical summand inclusions $E_j$. Summing over the
    sectors shows that the two adjacent placements of $Q$ commute.

    If $U$ is the common isometry, then transport gives the identity
    $(U\otimes U)Q(U^*\otimes U^*) =P_{G_2(B)}$,
    since the range on either side is
    $\bigvee_j G_2(A_j)=G_2(B)$. Hence it is the complement of the canonical
    two-site parent interaction of $B$. The two adjacent three-site parent
    terms therefore commute. Cyclic transport gives adjacent commutation on
    every chain with $N>2$; disjoint local terms commute by locality, which
    proves the claim.
\end{proof}

\begin{theorem}[Nearest-neighbor parent ground spaces for distinct multiplicity-one RFP sectors]%
    \label{thm:rfp_multiplicity_one_bnt_parent_ground_spaces}
    \lean{MPSTensor.IsBNTCanonicalForm.rfp_hasParentHamiltonianGroundSpaceSpanning_basisDirectSum}
    \leanok
    \uses{def:sector_bnt_cf, def:mps_transfer_idempotence,
        def:parent_hamiltonian_ground_space_spanning}
    Let $A_1,\ldots,A_g$ be the distinct basis tensors of a BNT canonical
    form, and let $B=\bigoplus_{j=1}^g A_j$ contain one copy of every basis
    tensor. If $B$ is a renormalization fixed point, then for every $N>2$,
    \begin{align}
        \ker H_2^{(N)}(B)
        &=\spn\bigl\{V^{(N)}(A_j):j=1,\ldots,g\bigr\}. \notag
    \end{align}
    This is the multiplicity-one distinct-sector case of
    \cite[Theorem~3.10]{Cirac2016MPDO_arXiv}. It does not include repeated
    copies of a basis tensor or arbitrary raw sector weights.
\end{theorem}

\begin{proof}\leanok
    The residual isometry equations imply that the simultaneous one-site
    matrices $(A_j^i)_j$ span the product of the block matrix algebras. Hence
    the same is true for every positive word length. The one-step restriction
    intersection therefore gives
    $\mathcal G_{N,2}(B) \subseteq \bigvee_j G_N(A_j)$.
    Comparing this block decomposition before and after moving the periodic
    cut by one site shows that every block boundary matrix commutes with the
    corresponding one-site matrices. Thus each block component lies in
    $\mathcal G_{N,2}(A_j)$, and
    $\mathcal G_{N,2}(B) =\bigvee_j\mathcal G_{N,2}(A_j)$.
    Each block is a normal renormalization fixed point, so one-site
    injectivity and the injective periodic-chain theorem identify its chain
    space with the line spanned by $V^{(N)}(A_j)$. The finite-volume kernel
    identity completes the proof.
\end{proof}

\begin{theorem}[Forward NNCPH ground spaces for distinct multiplicity-one RFP sectors]%
    \label{thm:rfp_multiplicity_one_bnt_has_nncph_ground_spaces}
    \lean{MPSTensor.IsBNTCanonicalForm.rfp_implies_hasNNCPHGroundSpaces_basisDirectSum}
    \leanok
    \uses{def:sector_bnt_cf, def:mps_transfer_idempotence,
        def:has_nncph_ground_spaces}
    Let $A_1,\ldots,A_g$ be the distinct basis tensors of a BNT canonical
    form, and let $B=\bigoplus_{j=1}^g A_j$ contain one copy of every basis
    tensor. If $B$ is a renormalization fixed point, then for every $N>2$ its
    translated two-site parent interactions commute and
    \begin{align}
        H_2^{(N)}(B)V^{(N)}(B)&=0, \notag\\
        \ker H_2^{(N)}(B)
        &=\spn\bigl\{V^{(N)}(A_j):j=1,\ldots,g\bigr\}. \notag
    \end{align}
    Hence $B$ satisfies the all-chain nearest-neighbor commuting
    parent-Hamiltonian ground-space condition. This is the multiplicity-one
    distinct-sector forward implication in
    \cite[Definition~3.9 and Theorem~3.10]{Cirac2016MPDO_arXiv}. It does not
    include repeated copies or arbitrary raw sector weights.
\end{theorem}

\begin{proof}\leanok
    \uses{thm:rfp_multiplicity_one_bnt_nncph,
        thm:rfp_multiplicity_one_bnt_parent_ground_spaces,
        rem:nncph_ground_spaces_decomposition}
    Theorem~\ref{thm:rfp_multiplicity_one_bnt_nncph} gives pairwise
    commutation of the translated two-site interactions for every $N>2$.
    The parent-Hamiltonian construction gives the zero-energy equation.
    Theorem~\ref{thm:rfp_multiplicity_one_bnt_parent_ground_spaces} identifies
    its kernel with the span of the periodic matrix product vectors of the
    distinct sectors. These are the three clauses of the all-chain condition.
\end{proof}
